\subsection{Spectral Weighting of Regional Variance}\label{sec:theory}
\label{sec:theory:bounds}\label{sec:mechanism}\label{sec:mechanism:loading}

Covariance allocation is a weighted sum of tangent energies. Its modal quotient
connects the scale curves to the trained features; coarse spectral bounds describe
what remains identifiable after grouping those modes.

\emph{Modal weights and limits.}
Let $F=V\Lambda V^\top$, with eigenvalues $\lambda_i\ge0$, largest eigenvalue
$\lambda_{\max}$ and orthonormal eigenvectors $v_i$, and write
$c_i(x)=g(x)^\top v_i$. At unit noise,
\begin{equation}\label{eq:sigma-sum}
 \sigma^2(g)=g^\top(\alpha I+F)^{-1}g
           =\sum_i\frac{c_i^2}{\lambda_i+\alpha}.
\end{equation}
For residual variance $v$, each term is $v c_i^2/(\lambda_i+\gamma)$,
$\gamma=\alpha v$. Defining $a_{r,i}=\langle c_i(x)^2\rangle_r$ gives
\begin{equation*}\tag{3a}\label{eq:modal-quotient}
 (\Rrms(\gamma))^2=
 \frac{\sum_i a_{\mathcal S,i}/(\lambda_i+\gamma)}
      {\sum_i a_{\mathcal B,i}/(\lambda_i+\gamma)}.
\end{equation*}
For these finite sums,
\[
 \lim_{\gamma\to\infty}\Rrms=\sqrt{A_{\mathcal S}/A_{\mathcal B}},
 \quad A_r=\langle\|g\|^2\rangle_r,
 \qquad
 \lim_{\gamma\downarrow0}\Rrms=\sqrt{K_{\mathcal S}/K_{\mathcal B}},
 \quad K_r=\langle\|\Pi_{\ker J}g\|^2\rangle_r,
\]
provided $A_{\mathcal B}>0$ and $K_{\mathcal B}>0$, respectively. Multiplication
by $\gamma$ gives both termwise limits. The QR-complement ratio equals the
whole-kernel ratio when $J$ has full row rank. For $A_{\mathcal S},A_{\mathcal B}>0$,
at strong precision, writing
$\mu_r=\sum_i a_{r,i}\lambda_i/A_r$ gives the more precise expansion
\begin{equation}\label{eq:r3-strong-expansion}
 (\Rrms(\gamma))^2=\frac{A_{\mathcal S}}{A_{\mathcal B}}
 \left[1-\frac{\mu_{\mathcal S}-\mu_{\mathcal B}}{\gamma}
       +O\!\left(\frac{\lambda_{\max}^2}{\gamma^2}\right)\right].
\end{equation}
This follows by expanding each $(\lambda_i+\gamma)^{-1}$ and dividing the two
sums. Neither endpoint bounds intermediate contrast, and monotonicity need not
hold. For example,
$\lambda=(0,1,100)$, $a_{\mathcal B}=(1,1,100)$ and
$a_{\mathcal S}=(1,100,1)$ give both limits equal to $1$, but
$\Rrms(1)=\sqrt{10304/503}\approx4.53$.
The measured curves and strong-precision checks are in
Appendix~\ref{app:r2:modal}. The conditional limits and counterexample are
Lean-verified (Table~\ref{tab:lean}); expansion~\eqref{eq:r3-strong-expansion}
remains outside verification.
Figure~\ref{fig:n1-scales} changes the modal weights while holding the trained
tangents fixed. The predictive centres share the covariance contrast but have different error targets.

\emph{Two-subspace bounds.} We group the modes by a spectral cutoff.
\begin{samepage}
\begin{definition}[Spectral split and regional loadings]\label{def:split}
At threshold $\epsilon>0$, let $\Pi_{\rm sl}$ project onto modes with
$\lambda_i\le\epsilon$ and $\Pi_{\rm st}$ onto those with $\lambda_i>\epsilon$.
Define
\begin{equation*}
 Q_{r,\rm sl}=\langle\|\Pi_{\rm sl}g\|^2\rangle_r,
 \qquad Q_{r,\rm st}=\langle\|\Pi_{\rm st}g\|^2\rangle_r.
\end{equation*}
\end{definition}
\end{samepage}
For an individual vector, the lower and upper envelopes are
\begin{equation}\label{eq:variance-envelopes}
 \begin{aligned}
 m_\epsilon(g)&=\frac{\|\Pisl g\|^2}{\alpha+\epsilon}
              +\frac{\|\Pist g\|^2}{\alpha+\lambda_{\max}},\\
 M_\epsilon(g)&=\frac{\|\Pisl g\|^2}{\alpha}
              +\frac{\|\Pist g\|^2}{\alpha+\epsilon}.
 \end{aligned}
\end{equation}

Theorem~\ref{thm:ceiling}, proved in Appendix~\ref{app:aggregation}, bounds each
pointwise variance between these envelopes. Averaging its termwise inequalities
gives the primary regional bound
\begin{equation}\label{eq:regional-rms-sandwich}
 L_{\rm reg}(\epsilon)\le(\Rrms)^2\le U_{\rm reg}(\epsilon).
\end{equation}
Here $L_{\rm reg}(\epsilon)=\langle m_\epsilon(g)\rangle_{\mathcal S}/\langle M_\epsilon(g)\rangle_{\mathcal B}$
and $U_{\rm reg}(\epsilon)=\langle M_\epsilon(g)\rangle_{\mathcal S}/\langle m_\epsilon(g)\rangle_{\mathcal B}$,
assuming positive total bulk energy. This is exact
finite-sum aggregation, without a Jensen step. For common variance $v$, replace
$\alpha$ by $\gamma$; the common factor cancels. The envelopes are sharp over
spectra compatible with these summaries (Proposition~\ref{prop:sharp}, proved in
Appendix~\ref{app:sharpness}), but they discard energy placement within each block.
They can therefore remain wide even after intersecting thresholds.
Appendix~\ref{app:aggregation} treats secondary aggregations; Appendix~\ref{oa:lean},
Table~\ref{tab:lean}, records the formal verification scope.
